% Pending full proof: A.139, 2026-11-19. No promotion before audit.
\begin{saresult}{A.139}{Proved under stated hypotheses}{The ordinary TOP concordance space of a sphere is the fibre of finite-rank stabilization}
For $n\ge2$, let
\[
 Y_n=\operatorname{hofib}\bigl(B\mathrm O(n)\longrightarrow
                                     B\mathrm{TOP}(n)\bigr),
 \qquad
 \mathrm{TOP}(n)=\operatorname{Homeo}(\mathbb R^n;0),
                                                        \tag{A.139.1}
\]
where the homeomorphism group has its actual compact-open topology.
For the ordinary bottom-fixed, free-top concordance space there is
an equivalence
\[
 P_{\mathrm{TOP}}(S^n)
   \simeq\operatorname{hofib}_*\bigl(Y_n\longrightarrow Y_{n+1}\bigr),
                                                        \tag{A.139.2}
\]
where the map is product with a line and all basepoints are standard.
In particular
\[
 \pi_i P_{\mathrm{TOP}}(S^4)=
 \pi_i P_{\mathrm{PL}}(S^4)=0\qquad(0\le i\le3).          \tag{A.139.3}
\]
For $i=0$ this means that the component set is a singleton.
No assertion about $\pi_4$ is made. The model retains the
finite-rank TOP groups; it is a sphere result and does not compute
the general remaining low-dimensional concordance spaces.
\end{saresult}
\begin{proof}
\paragraph{Conventions and primary inputs.}
Every topological group below is replaced by its singular simplicial
group. In particular orthogonal groups carry continuous singular
chains, and their inclusion in TOP is literal. Let $\mathcal G_n$
be the simplicial group of germs at $0$ of projection-preserving
homeomorphisms fixing the zero section on $\Delta^q\times U$, with one common open
neighbourhood $U$ for the entire simplex. Equality means equality
on a smaller common neighbourhood. These are the finite-parameter
germs of A.132, not block maps.

Salvatore--Turchin, Section~3.4, equations~(3.5)--(3.7),
uses precisely this finite-parameter germ convention
\cite{SalvatoreTurchin2026}. The globalization input is
Kuiper--Lashof, Lemma~3, printed p.~11, with proof on pp.~12--13
\cite{KuiperLashof1966I}. It extends a fibrewise embedding over a
parameter simplex, after shrinking its source, to a fibrewise sphere
homeomorphism preserving the two marked points. The relative
extension argument is given immediately after the lemma.
We do not use a smooth delooping theorem in dimension four.

\paragraph{Global and spherical representatives of germs.}
Restriction gives a degreewise surjective simplicial group map
\[
 T_n:=\operatorname{Sing}\operatorname{Homeo}(\mathbb R^n;0)
        \longrightarrow\mathcal G_n.                  \tag{A.139.5}
\]
Indeed compactness of each simplex gives a common source ball and
a common image chart for its germ representative. Fixed changes of
scale put it into the setting of Kuiper--Lashof Lemma3. The resulting
sphere extension fixes both $0$ and $\infty$; removing $\infty$
gives a global family on $\mathbb R^n$. A degreewise surjective
map of simplicial groups is a Kan fibration.

The kernel of (A.139.5) consists of families that are the identity on
one common ball about $0$. Inversion in the sphere interchanges
$0$ and $\infty$ and identifies this kernel with compactly
supported families on $\mathbb R^n$ that fix $0$. Such families
are weakly contractible by the Alexander formula
$a_s(x)=s\,a(x/s)$, $s>0$, with $a_0=1$.
Each finite family has common bounded support; this gives uniform
continuity of the map and inverse at $s=0$, and preserves the fixed
point. It follows that (A.139.5) is a weak equivalence, including its
component assertion.

Fix $p\in S^n$, choose an orthogonally equivariant chart at $p$,
and let $H_n=\operatorname{Sing}\operatorname{Homeo}(S^n;p)$.
The germ map
\[
                        H_n\longrightarrow\mathcal G_n
                                                               \tag{A.139.6}
\]
is likewise degreewise surjective by the same original lemma.
Its kernel is the group of sphere families that are the identity
on a common neighbourhood of $p$. On $S^n\setminus\{p\}$ it is
the group of compactly supported families. The same Alexander
formula contracts each finite family; there is no additional fixed
point condition in this instance. Hence (A.139.6) is also a weak
equivalence. These statements use common supports for finite
families; no claim about a union with a different topology is needed.

\paragraph{The first sphere fibre model, before identifying its map.}
Let $A_{n+1}=\operatorname{Sing}\operatorname{Homeo}(D^{n+1};0)$.
The germ restriction
$A_{n+1}\to\mathcal G_{n+1}$ is a Kan fibration. The proof in
A.132 is dimension-independent: the Edwards--Kirby local extension
operator is fixed on a neighbourhood of the identity; a common
source/image ball and a single finite mesh for each compact
parameter family extend successive increments, keeping prescribed
horns exactly fixed. No globalization of a disk vertex is needed
to prove this fibration assertion.

Its identity fibre is the filtered simplicial union of
$\operatorname{Sing}\operatorname{Homeo}(D^{n+1}\ \mathrm{rel}\
D_\epsilon^{n+1})$. Each stage is the ordinary free-outer-boundary
concordance space of $S^n$. The radial conjugation and compression
argument in A.132 applies in the same dimension-independent way:
any finite parameter family and any relative nullhomotopy admit a
common $\epsilon$, and inclusions of stages are weak equivalences.
Consequently the fibre is weakly equivalent to
$P_{\mathrm{TOP}}(S^n)$, with the inner face fixed and the outer
face free.

Boundary restriction $A_{n+1}\to
G_n:=\operatorname{Sing}\operatorname{Homeo}(S^n)$ is a weak
equivalence. It has the radial cone section, and its kernel consists
of boundary-fixed disk homeomorphisms also fixing $0$, contracted
by the relative Alexander formula. Extend every disk map outside
the disk by its coned boundary restriction. This gives a group map
$A_{n+1}\to T_{n+1}$ that commutes with germ restriction.
For the cone section its composite is the actual radial cone map
\[
 c:G_n\longrightarrow T_{n+1},\qquad
 c(h)(r\omega)=r h(\omega).                             \tag{A.139.7}
\]
Together with (A.139.5), these weak equivalences give
\[
 P_{\mathrm{TOP}}(S^n)\simeq\operatorname{hofib}_1(c).
                                                               \tag{A.139.8}
\]
The full germ equivalence (A.139.5) permits all components in this
formula; the proof is not restricted to an unproved identity-component
replacement of $G_n$.

\paragraph{Orthogonal quotients and the point stabilizer.}
The map (A.139.7) restricts to the standard linear inclusion on
$\mathrm O(n+1)$. Quotienting both its source and target by this
same subgroup therefore preserves its based homotopy fibre:
\[
 \operatorname{hofib}(c)\simeq
 \operatorname{hofib}\bigl(
 G_n/\mathrm O(n+1)\longrightarrow
 T_{n+1}/\mathrm O(n+1)\bigr).                          \tag{A.139.9}
\]
This follows from the map of principal orthogonal fibrations,
whose map on fibres is the identity. Quotients here can equivalently
be taken in simplicial groups or by homotopy orbits: the action is
free and the orthogonal principal bundles are locally trivial.

The inclusion of the point stabilizer induces an equivalence
\[
 H_n/\mathrm O(n)\longrightarrow G_n/\mathrm O(n+1).
                                                               \tag{A.139.10}
\]
For actual topological quotients it is a homeomorphism. Every right
coset has a representative fixing $p$, because $\mathrm O(n+1)$
acts transitively on $S^n$, and two such representatives differ
by exactly $\mathrm O(n)$. Local continuous choices of a rotation
taking $p$ to a nearby point provide the local inverses. Passing
to singular complexes gives (A.139.10); a principal orthogonal bundle
over a simplex has a section, so the simplicial and topological
quotient conventions agree up to the indicated equivalence.

Choose the chart at $p$ to be $\mathrm O(n)$-equivariant. Then
(A.139.6) is exactly the identity on that orthogonal subgroup and
gives a weak equivalence
\[
 H_n/\mathrm O(n)\xrightarrow{\simeq}
                \mathcal G_n/\mathrm O(n)\simeq Y_n.   \tag{A.139.11}
\]
Similarly (A.139.5) identifies the target quotient in (A.139.9) with
$Y_{n+1}$. It remains essential to identify the arrow between them.

\paragraph{Radial extension gives the actual stabilization map.}
Take $p=(0,\ldots,0,1)$ and $h\in\operatorname{Homeo}(S^n;p)$.
Its cone $c(h)$ fixes every point of the ray through $p$. Therefore
\[
 F_s(h)=\tau_{-sp}\,c(h)\,\tau_{sp},\qquad0\le s\le1,
                                                               \tag{A.139.12}
\]
is a continuous homotopy of group homomorphisms from $H_n$ to
$T_{n+1}$; every map fixes $0$. For $A\in\mathrm O(n)$,
embedded as $\operatorname{diag}(A,1)$, one has
$F_s(A)=\operatorname{diag}(A,1)$ for every $s$. Thus this homotopy
is exact on the orthogonal subgroup and descends to the quotients.
It moves the germ of $c(h)$ at $0$ to its germ at $p$.

Use the local chart
\[
 \theta(u)=(u,\sqrt{1-\|u\|^2}),\qquad
 \psi(u,t)=(1+t)\theta(u)-p.                            \tag{A.139.13}
\]
It is defined near $(u,t)=0$, has derivative the identity there,
and commutes with $\operatorname{diag}(A,1)$ for every
$A\in\mathrm O(n)$. Exact substitution gives
\[
 \psi^{-1}F_1(h)\psi(u,t)
                    =(\theta^{-1}h\theta(u),t).       \tag{A.139.14}
\]
The right-hand side is precisely product with a line on the
point germ of $h$.

Finally the coordinate change in (A.139.14) introduces no additional
map. The germs $\psi_r(z)=r^{-1}\psi(rz)$ for $r>0$, completed
at $r=0$ by $D\psi(0)=1$, form a continuous family of local
diffeomorphisms and their inverses on a common neighbourhood.
They commute with $\mathrm O(n)$ throughout. Conjugating by this
fixed family gives a homotopy, relative to that subgroup, from
the coordinate-conjugated germ in (A.139.14) to the original germ.
For any finite parameter family of $h$ a common smaller chart
works, by compactness. The same formulas preserve parameter faces.
Thus under (A.139.10)--(A.139.11) the map in (A.139.9) is the ordinary stabilization
$Y_n\to Y_{n+1}$ up to the specified based homotopy. This proves
the formula (A.139.2).

\paragraph{Safe low-degree consequence for the four-sphere.}
The corrected Freedman--Quinn statement, audited in A.132 using
Gabai--Gay--Hartman--Krushkal--Powell, Section~1.1, Theorem~1.3
\cite{GabaiGayHartmanKrushkalPowell2026}, says that
$Y_4\to\mathrm{TOP}/\mathrm O$ is5-connected. Hence it is an
isomorphism on $\pi_i$ for $i\le4$. The finite-rank statement
Kirby--Siebenmann, Essay~V, Theorem~5.5(I), p.~251
\cite{KirbySiebenmann1977}, gives
$\mathrm{TOP}_5/\mathrm{PL}_5\simeq K(\mathbb Z/2,3)$.
Together with the already audited
$\pi_i(\mathrm{PL}_5/\mathrm O_5)=0$ for $i<5$, this shows that
$Y_5\to\mathrm{TOP}/\mathrm O$ is an isomorphism for $i\le4$.
Naturality implies the same for $Y_4\to Y_5$.
The fibre sequence in (A.139.2) therefore gives
\[
                \pi_iP_{\mathrm{TOP}}(S^4)=0
                           \quad(0\le i\le3).         \tag{A.139.15}
\]
For PL the audited fibre formula A.128 and the corresponding
low-degree vanishing of $\mathrm{PL}_4/\mathrm O_4$ and
$\mathrm{PL}_5/\mathrm O_5$ give the same conclusion.
No vanishing of $\pi_4$ is inferred: that would require control
of $\pi_5Y_4\to\pi_5Y_5$, which the cited inputs do not supply.
No compatibility of an arbitrary model of PL/orthogonal groups
with continuous orthogonal singular chains is silently assumed.

\end{proof}
